%(Paragraph 1: Explain sequential recommendation.)
%Capturing users' preferences from their history is essential for making effective recommendations. Traditional methods, such as Collaborative Filtering (CF), focus on modeling the relationship between users and items based only on their interaction history without considering the temporal information. However, temporal information is essential because users' preferences and items' characteristics might evolve over time. Furthermore, users' preferences might depend heavily on the context (e.g. one is likely to be interested in buying a keyboard after purchasing a desktop). Therefore, sequential recommendation aims to model the sequences of users' past behaviors to predict the next set of items that they are likely to prefer.

Capturing users' preferences from their history is essential for making effective recommendations, because users' preferences and items' characteristics are both dynamic. Furthermore, users' preferences depend heavily on the context (e.g., one is likely to be interested in buying a keyboard after purchasing a desktop). Therefore, sequential recommendation aims to predict the next set of items that users are likely to prefer by exploiting their history.


%Paragraph 2: Brief touch on recent methods. Finish with SASRec, TiSASRec? and BERT4Rec.
%There have been many attempts to apply various algorithms to better understand the sequential behaviors of users. Works such as \cite{rendle2010factorizing, SIMILARITYMC, MCAULEYMC2016, MCAULEYMC2017, he2016fusing} employed Markov-Chain models by assuming that users' behaviors depend solely on their last N behaviors. Some tried to incorporate time information in the traditional factorization methods \cite{TIMESVD++}. More recent works adopted deep learning techniques such as RNN~\cite{GRU4Rec, hidasi2018recurrent, RNN4Rec, HIERARCHICALGRU4REC, DREAM, wu2017recurrent, wu2017sequential, DINEVOLUTION}, CNN~\cite{CASER, CNNGENERATIVE, HIERARCHICALCNN}, GNN~\cite{SESSIONGNN} for sequence modeling. Especially, the latest success came from self-attention mechanisms. For example, SASRec~\cite{SASRec} designed a strong model that runs entirely on self-attention blocks. As more recent work, BERT4Rec~\cite{BERT4Rec} improved SASRec~\cite{SASRec} by replacing their model with BERT~\cite{BERT}, which was originally a huge success in NLP area.

%Paragraph 3: Point out the limitations of recent methods.
%Many algorithms have been tried to better understand the sequential behaviors of users~\cite{SASRec, BERT4Rec, hidasi2018recurrent, CASER, M3, rendle2010factorizing, DIN, MCAULEYMC2017}. Although these studies showed excellent performance, most of them use time information only to order the user's behaviors into a sequence. This can be considered a loss since we can embed time information to extract more patterns from the sequence. While recent works such as TiSASRec~\cite{TISASRec} successfully incorporated time information into their recommendation model, their usage of time was limited to a single embedding scheme. Assuming that adding novel positional embedding benefits performance, it is natural to investigate more forms of temporal encodings and a model architecture that can fully leverage such diversity.


Many algorithms have been proposed to better understand the sequential history of users~\cite{SASRec, BERT4Rec, hidasi2018recurrent, CASER, M3, rendle2010factorizing, DIN, TRANSREC}. Despite their excellent performance, most of them ignore the interactions' timestamp values. While recent works such as TiSASRec~\cite{TISASRec} successfully incorporated time information, their usage of time was also limited to a single embedding scheme. Since the timestamp values are rife with information, it would be beneficial to explore many forms of temporal embeddings that can fully extract diverse patterns, and model architectures that can fully leverage such diversity.



%Paragraph 4: Our solution.
%In this paper we propose \textbf{MEANTIME}(\textbf{M}ixtur\textbf{E} of \textbf{A}tte\textbf{NTI}on mechanisms with \textbf{M}ulti-temporal \textbf{E}mbeddings) to address the issues mentioned above.Recently, CTA\cite{CTA} demonstrated that using multiple kernels for modeling temporal influence can benefit performance. Inspired by this, \textbf{MEANTIME} introduces multiple temporal embeddings to better encode both absolute and relative positions of user-item interactions in a sequence. Moreover, we employ multiple self-attention blocks simultaneously to handle each positional encoding separately. This is profitable because each block can function as a sort of expert that focuses on a particular pattern of user behaviors. For example, one block could concentrate on short-term effects of nearby interactions, while another block could extract long-term effects from distant interactions. Some other blocks could disregard the positional bias, thereby only examining the relationship between items. Extensive experiments show that \textbf{MEANTIME} outperforms several state-of-the-art sequential recommendation algorithms on many real-world datasets.\\

In this paper we propose \textbf{MEANTIME} (\textbf{M}ixtur\textbf{E} of \textbf{A}tte\textbf{NTI}on mechanisms with \textbf{M}ulti-temporal \textbf{E}mbeddings) which introduces multiple temporal embeddings to better encode both absolute and relative positions of a user-item interaction in a sequence. Moreover, we employ multiple self-attention heads that handle each positional embedding separately. This is profitable because each head can play a role of expert that focuses on a particular pattern of user behaviors. 
%For example, one block could concentrate on short-term effects of nearby interactions, while another block could extract long-term effects from distant interactions. Some other blocks could disregard the positional bias, thereby only examining the relationship between items. 
Extensive experiments show that \textbf{MEANTIME} outperforms state-of-the-art algorithms on real-world datasets. Our contributions are as follows:
\begin{itemize}
\item We propose to diversify the embedding schemes that are used to encode the positions of user-item interactions in a sequence. This is done by applying multiple kernels to positions/values of timestamps in a sequence to create unique embedding matrices.
\item We also propose a novel model architecture that operates multiple self-attention heads simultaneously, where each head specializes in extracting certain patterns from the users' behaviors by utilizing one of the positional embeddings we propose above.
\item We conducted thorough experiments to show the effectiveness of our method on real-world datasets. We also provide a comprehensive ablation study that helps understand the effect of various components in our model.
\end{itemize}
